\section{The Purpose of Life Revealed}

I first wrote about the purpose of life several years ago in Not Finding Life's Purpose\footnote{\url{https://gornahoor.net/?p=38}}. We revisited the topic a few months ago, but this time we want to dig deeper by contemplating the \emph{meaning} of life, as the recent posts about Keyserling emphasize.

How you answer this will depend on your idea of what it means to be a human being. Of course, the better idea will account for more facts. On the first level, we can take the position that a human being is simply an animal and is therefore subject to biological laws. Since the primary biological law is the preservation of the species, then man's primary purpose is reproduction. But this is not his meaning, that is, it is not self-possession, because it depends on an outside force such as the mythical ``selfish gene'' or strong sexual urges that lead to reproduction in the natural world.

\paragraph{Natural End}


Dr. Bruce Charlton points out in Longevity versus real health\footnote{\url{http://charltonteaching.blogspot.com/2014/08/longevity-versus-real-health-more.html}} that the prolongation of human life beyond the prime reproduction years is abnormal or even pathological. From that perspective, it makes no biological sense. Then why, unlike physics, is biology not subject to the conservation of energy through following the path of least resistance? From a strictly biological point of view, this prolongation of life is just a waste of energy with no real explanation. Analogous phenomena are not seen in physics. What makes sense, from this point of view, is the situation of the drone honey bee who dies immediately after coitus.

On the other, we know that the human person cannot be fully understood solely in biological terms. Rather, our biological nature serves as the ``not I'', the substrate, the field of activity where we can develop our higher nature. In a misunderstanding of Thomism, the West for some time held to a two tier solution: we have both a natural and a supernatural purpose or end. That point of view is hard to maintain, since the two ends hardly seem reconcilable. In practice, it led to the diminution of the supernatural end, which is either forgotten or left undeveloped, so that most people simply pursue a natural end.

I certainly don't want to imply that this is a new phenomenon, since an ancient text like the \emph{Bhagavad Gita} says essentially the same thing. The real point is that the pursuit of solely natural ends represents a deformation of human life, not a sign of health. Dr. Charlton's ideal of a biologically healthy society would include

\begin{itemize}


\item early death

\item the elimination of the weak and unfit

\item high child mortality rate.
\end{itemize}


Those goals may be true strictly from a eugenic point of view, but it presumes that a longer life is dysgenic. There are two impulses at play, then, in the human condition, not necessarily compatible: the desire for longevity of the individual and the desire for posterity. Rene Guenon gives us an esoteric meaning for those desires, so we turn to that next.

\paragraph{Exoteric Salvation}


In Taoism, longevity, one of the \emph{Four Happinesses}, means everlasting life. Obviously, longevity is a happiness only if that life occurs in a state of bliss, which we call Heaven, and not in a Hell. Hence, salvation is the achievement of a postmortem state in Heaven. Biological and social life can therefore only be the field in which such a salvation can be achieved. This state is the restoration of the primordial state, which was lost somehow (this is not the place to get into this). Since the subtle rules the dense, this loss adversely affected life on the biological plane. Apparently, this immortal state occurs only on a subtle level, since the will to power to bring it down to the physical and biological level is not strong enough, although that possibility is still left open at some future time.

Salvation is experienced as something gratuitous, something that lifts us out of our ordinary life to an awareness of transcendence. This is illustrated, for example, in the document \emph{Lumen Gentium} of the Roman Church:

\begin{quotex}
their exalted position results, not from their own merits, but from the grace of Christ. If they fail to respond in thought, word and deed to that grace, not only shall they not be saved, but they shall be the more severely judged.
\end{quotex}


There we see the two points. First, the experience arrives from a transcendent source, unexpectedly and undeserved. Second, our life is seen as a battlefield in which spiritual warfare is to be carried out. The Reformation distorted this teaching in an unbalanced way, admitting only the first point, but not the second. Because of the false doctrine of ``once saved, always saved'', there is no inner battle necessary.

Although I don't want to initiate or engage in a dispute over exoteric doctrines, I should point out that the Roman church teaches that this transcendent grace may also be available to those following other traditions. I only bring this up since some have accused me of making scandalous innovations in this regard.

For most people, salvation in this sense is sufficient. Let's do a phenomenological analysis of the salvation experience. First of all, there is the experience of a new meaning to life, a meaning experienced on a deeper level. Biological life can no longer be understood as sufficient in itself, but our life is experienced as open to transcendence. Nevertheless, there are some features of full personhood that are missing.

Firstly, the experience is ``spontaneous''. It arrives into our consciousness as fully formed. The person is an ``I'', but in this case, the I is only passive to that experience. The other characteristic of the person is self-possession. The person must own the experience in an active way. This leads us to a still deeper understanding of the meaning of life.

\paragraph{Esoteric Path}


Some few will be motivated to understand salvation in an active way. This involves moving well beyond biological life and the spontaneity of exoteric ideas of salvation. The person is the center of consciousness and will, so he is led to understand his own states of consciousness, to find his own life. As Aristotle wrote, ``The soul is all that it knows.'' Hence, the salvation of the soul is tied to a certain kind of knowing. This is called ``gnosis'', and is an active knowing, not the passive reception of intellectual knowledge.

\textbf{St Teresa of Avila} wrote that most errors in the spiritual life arise from a poor understanding of the nature of the soul, so we are on solid ground here. Our full salvation depends on our knowledge and understanding of our own souls. Since this is not a spontaneous process, it requires considerable efforts both to pay attention to our inner states and to engage in the spiritual battle in full consciousness. These efforts cannot be understood as ``works'' in the Biblical sense, as good things we do in order to be rewarded. Through such efforts, we may come to the realization of our true Self or real I. Or, as we have expressed it many times previously, this is the second birth of the Logos in our souls.

This is not the place to discuss all the details of this path, but rather it serves as a reminder of our true purpose in life. But to close, we now can better understand the point of a long life, despite its biological dysfunction. Salvation understood exoterically can easily devolve into mere repetition. The same spiritual battles are fought over and over, with the final hope of dying in a ``state of grace''. Esoterically understood, salvation is a process requiring self-knowledge and efforts.

\emph{A long life provides more opportunities for development and self-realization.}


\flrightit{Posted on 2014-08-07 by Cologero}

\begin{center}* * *\end{center}

\begin{footnotesize}
\begin{sffamily}
\texttt{Pickman on 2014-09-12 at 21:36 said:}

``Some few will be motivated to understand salvation in an active way. This involves moving well beyond biological life and the spontaneity of exoteric ideas of salvation. The person is the center of consciousness and will, so he is led to understand his own states of consciousness, to find his own life. As Aristotle wrote, ``The soul is all that it knows.'' Hence, the salvation of the soul is tied to a certain kind of knowing. This is called ``gnosis'', and is an active knowing, not the passive reception of intellectual knowledge.

So our full salvation depends on our knowledge and understanding of our own souls. Since this is not a spontaneous process, it requires considerable efforts both to pay attention to our inner states and to engage in the spiritual battle in full consciousness. These efforts cannot be understood as ``works'' in the Biblical sense, as good things we do in order to be rewarded. Through such efforts, we may come to the realization of our true Self or real I. Or, as we have expressed it many times previously, this is the second birth of the Logos in our minds.

Esoterically understood, salvation is a process requiring self-knowledge and efforts. A long life provides more opportunities for development and self-realization.''

What are these efforts: what is the manual? I'm going to put it in as blunt as terms as that.

There is a lot of poetic esoteric language, euphemisms, analogies and promises but little about the practical processes which is what most spiritual aspirants are hungry for. I would humbly request the author provide some direction/information in that vein.

\hfill

\texttt{Cologero on 2014-09-14 at 23:40 said:}

Well, Pickman, the first step is the realization that there is something wrong with ordinary life, so a man may become a ``spiritual aspirant''. Then he must learn true doctrine, which Gornahoor tries to provide. If the information is incomplete, that is because we have pointed to many authentic teachings, which is where the search must begin. Beyond that, one must find a ``way'' as you pointed out. For a few, the available exoteric path is insufficient to them, so they seek an esoteric path. Maybe you have found such a path or have considered being a monk. For those who live in the world, something else is needed.

For those who are not satisfied with an exoteric way, or cannot find a suitable esoteric way, we try to provide a little push so such a way can be recognized when it appears. We have certainly provided the means to a ``practical process'' for those willing to make the efforts. There is a \textit{discussion list}\footnote{\url{http://www.meditationsonthetarot.com/medtarot-discussion-list}} for those interested in Meditations on the Tarot.

More recently, we posted information on a personal study group on \textit{Achieving Gnosis in Practice}\footnote{\url{http://www.gornahoor.net/?p=7205}}.

It is at the top of the home page and hard to miss. If you are interested in participating, use the
\textit{``contact us'' form}\footnote{\url{http://www.gornahoor.net/?page_id=4143}} for more information.

\hfill

\texttt{Pickman on 2014-09-15 at 17:57 said:}

The participation group has been noted along with the highly informative articles provided. I will be in contact regarding further actions soon. I hope my inquiries may serve as a useful guide to others in what can often be a tortuous route, thanks for your reply.


\hfill

\end{sffamily}
\end{footnotesize}

